\section{局限、风险与合规}
\label{sec:limits}

本章集中说明系统的已知局限。我们把局限分为四类：评测证据本身的局限、计算口径的简化、工程层面的约束，以及使用风险与合规。前三类决定了本文结论能推广到多远，第四类决定了系统可以被怎样使用。

\subsection{评测证据的局限}

\paragraph{样本规模小。}真实年报的抽取核验只覆盖少数几家公司；研报质量评分只在两到三家公司上做过迭代对比；评审校准的劣化实验每家公司只有一组算子。这些数字足以说明机制“能工作”，不足以给出带置信区间的性能估计。团队正在按 \file{docs/eval-dataset-collection.md} 的清单收集约 15 家公司的年报、季报与卖方研报，并人工埋错构造纠错评测集。

\paragraph{版本对比不是对照实验。}研报得分从 v1 到 v2 的提升，同时包含了流水线改动（先搜集、深度专题）、评分标准修订与模型随机性的影响，我们没有在同一评分标准下对每项改动做单独的对照，也没有多次重复生成取均值。第 \ref{sec:exp} 节中的消融结果应被视为初步证据。

\paragraph{样例数据是虚构的。}财报分析、估值建模与纠错核查的回归评测使用虚构公司“星河智造”的样例年报与研报草稿，样例内部完全自洽，版式也比真实年报规整。样例指标衡量的是“引擎在理想输入上是否正确”，真实年报上的表现会明显下降，这也是我们单独报告真实年报抽取结果的原因。

\paragraph{评审与生成模型同族。}盲评质检员与生成角色都来自同一模型家族（分别为 Sonnet 与 Opus）。已有研究表明评审模型可能偏好与自身风格相近的文本\cite{panickssery2024self}。我们用信息隔离（评审看不到生成过程与模型名）、程序判定的 L0/L1 层与人工抽查部分缓解这一问题，但没有用不同家族的评审模型做交叉验证。

\paragraph{单次盲评。}正式得分来自单次盲评。重复评审的逐条一致率在 86\%～91\% 之间、总分极差 7～8 分，这意味着相差几分的两份研报，其优劣不能仅凭一次评分判断。

\subsection{计算口径的简化}

以下简化均已在代码与报告中注明，列在这里以便读者判断估值结论的适用范围：

\begin{itemize}
  \item \textbf{EBIT 口径。}EBIT 定义为利润总额 + 财务费用 − 投资收益 − 公允价值变动收益；资产处置收益、其他收益、减值与营业外收支仍留在 EBIT 内。剔除的投资类资产需要在假设文件的 \code{non\_operating\_assets} 中加回，未设置时估值偏保守，报告会给出提示。
  \item \textbf{经营费用是轧差项。}三表模型中的经营费用由毛利 − EBIT − 折旧摊销倒推，以保证历史期 EBIT 与 DCF 同口径；代价是经营费用不能按费用类别分别预测。
  \item \textbf{期初留存收益置零。}三表模型中权益合并列示，期初留存收益置 0，只影响列示，不影响勾稽。
  \item \textbf{流动性压力测试的符号。}压力测试沿用原插件口径 FCF $\times(1+\text{冲击})$，当 FCF 为负时冲击方向相反。为保持与参照实现一致我们未修改，报告中已注明。
  \item \textbf{金融企业。}FCFF 框架对银行、保险等金融企业并不适用，其资产负债表本身就是经营资产。对这类公司，系统的估值输出应只作参考。
\end{itemize}

\subsection{工程约束}

\begin{itemize}
  \item \textbf{不做 OCR。}扫描页直接拒收。部分上市公司的年报附件或早期年报为扫描件，这些文档目前无法处理。
  \item \textbf{季报表头。}规则通道不识别季报“本报告期 / 年初至报告期末”表头，会把它们当成年度数据，因此季报只接受模型通道抽取，规则通道不参与比对，双通道交叉校验在季报上失效。
  \item \textbf{规则通道覆盖。}股本变动表、政府补助口径以及部分扣非指标的季度别名，规则通道尚未覆盖。
  \item \textbf{成本与时延。}一份深度研报需要五个长会话加盲评，单次运行耗时以小时计，费用远高于一次 API 调用。这一成本对单家公司深度研究是可接受的，对大范围覆盖则不可接受。
  \item \textbf{外部来源的可靠性与时效。}来源登记保证“数字在摘句中”，不保证摘句所在网页本身准确，也不保证网页在报告发布后仍然可访问。系统检查来源日期以避免使用报告日之后的信息，但无法排除来源本身转载有误。
  \item \textbf{模型输出不可离线复现。}已保存输入上的程序计算可以逐位复算；重新生成文稿或重新盲评依赖模型环境，结果会有差异。我们保存了每个会话的完整事件流，但不宣称能复现模型输出本身。
\end{itemize}

\subsection{使用风险与合规}

\paragraph{不构成投资建议。}系统输出的评级由程序按基准情景每股价值相对参考股价的上行空间机械给出，依赖假设文件中的每一个参数。每份报告都写明假设、口径简化、适用范围与风险，并附“不构成投资建议”声明。系统不应被用作投资决策的唯一依据。

\paragraph{看起来专业的错误。}系统能生成排版精良、图表丰富、数字可追溯的研报，这本身带来一种风险：读者可能因为形式上的严谨而高估观点部分的可靠性。可追溯性只覆盖数字，投资主线、驱动判断与预期差分析仍是模型的观点。报告按“事实 → 推论 → 观点”分层，并在附录写出全文数字的出处构成，就是为了让读者看清哪些内容有程序保证、哪些没有。

\paragraph{版权与数据使用。}年报与公告是公开披露文件；卖方研报与付费数据库内容受版权保护。仓库只存储我们从公开文件中抽取的结构化结果与自己生成的文本，评测用的卖方研报只读不入库。外部来源登记只保存短摘句与链接。

\paragraph{第三方代码。}计算引擎的若干模块参照了一个开源插件的计算口径并以自研实现替代，参照关系、版本与许可证如实列于 \file{THIRD\_PARTY.md}（附录 \ref{app:thirdparty}），不冒充自研。
